% ======================================================================
\begingroup\sloppy
\section{Large-distance behaviour and integrability of the finite terms}\label{app:tails}
% ======================================================================
% Adapted from the verified Sections 6.3-6.7 of the previous draft (paper/draft/sec_finite.tex;
% notes paper/notes/fix_soft.tex, fix_tail.tex, ref3_finiteness.tex, ref3_prescriptions.tex).

This appendix gives the reasons for the prescriptions of Section~\ref{sec:fin_together}: which terms of the finite part have non-integrable tails at large transverse distance, how these tails cancel, and in which sense the remaining integrals converge. The special functions introduced here are not needed to state the result; they are used only to exhibit its large-distance behaviour. We use the abbreviations
\begin{align}\label{eq:Cf}
\mathcal C_{\rm f}\equiv\frac{1}{(2\pi)^3}\,\frac{g^4}{16\pi^5}\,\frac{1}{k^+}=\frac{\mathcal C}{\pi}\,,\qquad P^i\equiv\Kv^i(x'-w')\,,
\end{align}
with $\mathcal C$ of Eq.~\eqref{eq:dg_not}. A term is called integrable if its singularities at all coincidence points are integrable and if it decays faster than $1/\v z^2$ at large transverse position $\v z$ of the integrated gluon (in practice like $\log|\v z|/|\v z|^3$). A term $\propto e^{-i\v k\cdot\v z}/\v z^2$, in which the measured momentum flows through $\v z$, converges as a Fourier integral, as $\int_we^{-i\v k\cdot\v w}\Kv(x-w)$ does at LO, and is not counted as a tail.

\subsection{Closed forms of the longitudinal integrals}\label{app:closedforms}

When only the recoil phase depends on $\xi$, the integrals of Eqs.~\eqref{eq:plus_xi} and~\eqref{eq:inst_xi} are, with a real variable $s$ (${\rm Ci}$, ${\rm Si}$, $E_1$ the cosine, sine and exponential integrals),
\begin{align}\label{eq:recoilfns}
\int_0^1d\xi\,\Big[\frac1\xi\Big]_+e^{-is\xi}&=\int_0^1\frac{d\xi}{\xi}\big(e^{-is\xi}-1\big)\equiv\mathcal E(s)={\rm Ci}(|s|)-\gamma_E-\log|s|-i\,{\rm Si}(s)\,,\notag\\
\int_0^\infty d\xi\,\Big[\frac{1}{\xi(1+\xi)}\Big]_+e^{-is\xi}&=\int_0^\infty\frac{d\xi}{\xi(1+\xi)}\big(e^{-is\xi}-1\big)\equiv\Psi(s)=\mathcal E_\infty(s)-\widehat{\mathcal E}(s)\,,\notag\\
\int_0^\infty\frac{d\xi}{1+\xi}\,e^{-is\xi}&\equiv\widehat{\mathcal E}(s)=e^{is}E_1(is)\,,\qquad
\mathcal E_\infty(s)\equiv-\gamma_E-\log|s|-\frac{i\pi}{2}\,{\rm sgn}(s)\,,\notag\\
\int_0^1d\xi\,e^{-is\xi}&=\frac{1-e^{-is}}{is}\,.
\end{align}
For $|s|\to\infty$, $\mathcal E(s)=\mathcal E_\infty(s)-\frac{e^{-is}}{is}+O(s^{-2})$, $\Psi(s)=\mathcal E_\infty(s)-\frac{1}{is}+O(s^{-2})$ and $\widehat{\mathcal E}(s)\simeq\frac{1}{is}$; for $s\to0$, $\mathcal E(s)\simeq-is$, $\Psi(s)=O(s\log|s|)$ and $\widehat{\mathcal E}(s)\simeq\mathcal E_\infty(s)$. All these functions satisfy $f(-s)=f(s)^*$. $\mathcal E$ belongs to Group~I, where the $\xi$ integral runs up to one, and $\Psi$ to Group~III, where it runs to infinity; at large $|s|$ both tend to the same function $\mathcal E_\infty$, including its odd part.

\subsection{Soft remainders of the plus prescriptions}\label{sec:fin_soft}

\paragraph{Origin.} In Rows~I.III and~I.V the measured gluon in the amplitude, and in Rows~III.III, III.V and~III.VI in the conjugate amplitude, is part of a two-gluon state: a hard gluon at $\v w$ (resp.\ $\v w'$) and a gluon with $p^+=\xi k^+$ at $\v z$. The measured gluon is at the $+$-momentum weighted centre, so that the phase is $e^{-i\v k\cdot\v w}e^{-i\xi\v k\cdot(\v z-\v w)}$ (Group~I), resp.\ $e^{i\v k\cdot\v w'}e^{-i\xi\v k\cdot(\v z-\v w')}$ (Group~III); in both cases the exponent is $-i\xi s$ with $s=\v k\cdot(\v z-\text{hard position})$. The longitudinal integral of a term $\frac{d\xi}{\xi}g(\xi)$ of a Group~I row is split as
\begin{align}\label{eq:softsplit}
\int_{\Lambda/k^+}^{1}\frac{d\xi}{\xi}\,e^{-i\xi s}g(\xi)&=\log\frac{k^+}{\Lambda}\,g(0)+g(0)\,\mathcal E(s)+\int_0^{1}\frac{d\xi}{\xi}\,e^{-i\xi s}\big[g(\xi)-g(0)\big]+O\Big(\frac{\Lambda}{k^+}\Big)\,,
\end{align}
and a term $\frac{d\xi}{\xi(1+\xi)}g(\xi)$ of a Group~III row in the same way, with $\Psi(s)$ in place of $\mathcal E(s)$. Here $g(0)$ is the coefficient of the row in the JIMWLK list, which is kept exactly as in Section~\ref{sec:JIMWLK}; the last two terms together are the plus-prescription term of Eq.~\eqref{eq:plus_xi} with $f(\xi)=e^{-i\xi s}g(\xi)$. We call $g(0)\,\mathcal E(s)$, resp.\ $g(0)\,\Psi(s)$, the soft remainder of the term; the last term of Eq.~\eqref{eq:softsplit} is integrable by itself. Since $|\mathcal E(s)|,|\Psi(s)|\simeq\log|s|$ for $|s|\gg1$, the recoil cuts the soft integral off at $\xi\simeq1/|\v k\cdot(\v z-\v w)|$, i.e.\ at $p^+/p_\perp\simeq k^+/|\v k|$ for $p_\perp\sim|\v z-\v w|^{-1}$: the soft gluon is restricted to rapidities below that of the measured gluon.

\paragraph{The double logarithms.} Each soft remainder behaves at large $\v z$ like $\frac{T}{\v z^2}\mathcal E_\infty(\v k\cdot\v z)$, where $T$ is the $1/\v z^2$ tail of $g(0)$, and is not integrable by itself. Before the split~\eqref{eq:softsplit}, i.e.\ at finite $\Lambda$, every such row contains $\pi\log^2(k^+/\Lambda)\,T_{\rm row}+O(\log(k^+/\Lambda))$, from the region in which the soft gluon is far from the measured one and is cut off only by the recoil; in the rows without recoil the tails multiply single logarithms only. The coefficient of the double logarithm in the sum of all rows is $\sum_{\rm rows}T_{\rm row}$. This sum vanishes for an arbitrary projectile, in the symmetrised two-$\rho$ and three-$\rho$ sectors (numerically $|\sum T|<10^{-7}$ against single tails of order one to six, for $SU(2)$ and $SU(3)$; changing any one weight gives a sum of order one), and hence the double logarithms $\log^2(k^+/\Lambda)$ cancel in the cross section. What is left at large distance in the logarithmic terms is only the infrared behaviour of the JIMWLK term for a non-singlet pair of charges; it vanishes for a colour-diagonal correlator, Eq.~\eqref{eq:singlet}, and is not affected by the finite terms. The vanishing of $\sum T$ is explained by the pair structure.

\paragraph{The pairs.} The soft remainders cancel in pairs. Each Group~I term (soft gluon absorbed by the measured gluon in the amplitude: virtual, recoil on $\v w$) has one Group~III partner (soft gluon emitted by the measured gluon in the conjugate amplitude: real, recoil on $\v w'$) with the same coefficient, the same colour structure and the same kernel, except that $\v w$ is replaced by $\v w'$. With the colour operators of the Group~I members in their ordered form,
\begin{align}\label{eq:CABCD}
\mathcal L^a&\equiv\big[U^{ab'}(x')-U^{ab'}(w')\big]\rho^{b'}(x')\,,\notag\\
C_A&\equiv f^{abd}\,\mathcal L^a\,U^{de}(x)\big[U^{bc}(w)-U^{bc}(y)\big]\rho^e(x)\rho^c(y)\,,\notag\\
C_D&\equiv f^{abd}\,\mathcal L^a\,U^{de}(x)\big[U^{bc}(z)-U^{bc}(y)\big]\rho^e(x)\rho^c(y)\,,\notag\\
C_B&\equiv\mathcal L^a\Big[f^{adc}\,U^{db}(z)\,U^{ce}(w)\,\{\rho^e(y),\rho^b(x)\}-f^{abc}\,U^{cd}(y)\,U^{be}(x)\,\{\rho^d(y),\rho^e(x)\}\Big]\,,\notag\\
C_C&\equiv\big[U^{ab'}(x')-U^{ab'}(w')\big]\Big[N_c\,U^{ah}(x)-f^{adc}f^{hbe}\,U^{db}(z)\,U^{ce}(w)\Big]\rho^{b'}(x')\,\rho^h(x)\,,
\end{align}
and with the variables of Row~I.III (soft gluon at $\v z$, hard gluon at $\v w$, charges at $\v x'$, $\v x$, $\v y$; the terms of Row~I.V relabelled accordingly, and in Group~III $\v y'$ renamed $\v w'$), each pair has the form
\begin{align}\label{eq:pair}
\text{pair }X:\qquad\mathcal C_{\rm f}\,\alpha_X\,V_X\,C_X\Big[\kappa_X(w)\,\mathcal E\big(\v k\cdot(\v z-\v w)\big)-\kappa_X(w')\,\Psi\big(\v k\cdot(\v z-\v w')\big)\Big]
\end{align}
under $\int_{w,w'}e^{-ik\cdot(w-w')}\int_{x,y,x',z}$, with the entries of Table~\ref{tab:pairs}.
\begin{table}[t]
\centering
\renewcommand{\arraystretch}{1.3}
\footnotesize\setlength{\tabcolsep}{3.5pt}
\begin{tabular}{@{}c@{\hspace{5pt}}l@{\hspace{5pt}}l@{\hspace{5pt}}c@{\hspace{5pt}}c@{\hspace{5pt}}c@{\hspace{5pt}}l@{}}
\toprule
pair & Group I & Group III & $\alpha_X$ & $V_X$ & $\kappa_X(u)$ & tail $T_X$\\
\midrule
A & I.III, $\mathfrak K^{(w)}_{\rm I.III}$ & III.V & $-2i$ & $P\cdot\Kv(y-w)$ & $K(x,u,z)$ & $-2i\,P\cdot\Kv(y-w)\,C_A$\\
B & I.V, cubic & III.III, cubic & $-i$ & $P\cdot\Kv(y-w)$ & $K(x,u,z)$ & $-i\,P\cdot\Kv(y-w)\,C_B$\\
C & I.V, quad & III.III, quad & $-1$ & $P\cdot\Kv(x-w)$ & $K(x,u,z)-2K(w,u,z)$ & $+P\cdot\Kv(x-w)\,C_C$\\
D & I.III, $\mathfrak K^{(z)}_{\rm I.III}$ (last) & III.VI & $+2i$ & $P\cdot\Kv(x-w)$ & $K(y,u,z)$ & $+2i\,P\cdot\Kv(x-w)\,C_D$\\
\bottomrule
\end{tabular}
\caption{The four pairs of soft remainders, Eq.~\eqref{eq:pair}. $T_X\equiv\lim_{|z|\to\infty}\v z^2\alpha_XV_X\kappa_X(w)C_X$ is the $1/\v z^2$ tail of the Group~I member; that of the Group~III member is $-T_X$. ``Cubic'' and ``quad'' denote the cubic and the genuine two-$\rho$ part of a row; for pair~D the Group~I member is the last term of $\mathfrak K^{(z)}_{\rm I.III}$. In pair C, $C_C$ contains $U(\v w)$ also in its second term, and $K(w,u,z)$ is $1/(\v z-\v w)^2$ at $u=w$ and $K(w,w',z)$ at $u=w'$. In pair B the two recoil terms of Row~I.V are combined after the relabelling $\v y\leftrightarrow\v y'$ of the integrated charge positions, and in pair D Row~III.VI is written with $\v x\leftrightarrow\v y$.}
\label{tab:pairs}
\end{table}
That the Group~III member has the colour structure of its Group~I partner follows from $U^{d'b}(z)U^{cb}(z)=\delta^{d'c}$, $f^{c'd'a}f^{dd'a}=N_c\delta^{c'd}$ and a relabelling of colour indices only; no charges are reordered, so these are operator identities, valid for any $U$ at any $\v z$, and they hold separately for the symmetrised and for the contact terms, since the two members of a pair have the same ordered colour operator and hence the same decomposition~\eqref{eq:symord}--\eqref{eq:symr21}. They were also verified numerically, together with the statement that each entry of Table~\ref{tab:pairs} equals the row's own coefficient, in the two- and three-$\rho$ sectors for $SU(2)$ and $SU(3)$ (to $2\times10^{-15}$).

\paragraph{The two recoil sides combine without a condition on the target.} Each pair can be written as
\begin{align}\label{eq:rmv}
&\kappa(w)\,\mathcal E\big(\v k\cdot(\v z-\v w)\big)-\kappa(w')\,\Psi\big(\v k\cdot(\v z-\v w')\big)\notag\\
&\quad=\big[\kappa(w)-\kappa(w')\big]\mathcal E\big(\v k\cdot(\v z-\v w)\big)
+\kappa(w')\Big[\mathcal E\big(\v k\cdot(\v z-\v w)\big)-\Psi\big(\v k\cdot(\v z-\v w')\big)\Big]\,.
\end{align}
The first term is a real-minus-virtual combination times a kinematic logarithm: $\kappa(w)-\kappa(w')=O(|\v w-\v w'|/|\v z|^3)$ while $|\mathcal E|\le\log|\v z|+{\rm const}$; it has the structure of the remainder found in forward production~\cite{Altinoluk:2014eka} and is integrable. The second term is the mismatch of the recoil on $\v w$ and on $\v w'$: $\mathcal E(\v k\cdot(\v z-\v w))-\Psi(\v k\cdot(\v z-\v w'))=\log\frac{|\v k\cdot(\v z-\v w')|}{|\v k\cdot(\v z-\v w)|}+O(1/|\v k\cdot\v z|)$; the odd parts $-\frac{i\pi}{2}{\rm sgn}$ are the same on the two sides and cancel, and the remainder is $O(1/|\v z|)$ outside the strip $|\v k\cdot\v z|\lesssim|\v k||\v w-\v w'|$ and bounded by logarithms inside it, which gives a convergent integral. The colour structure $C_X$, with its whole dependence on $U(\v z)$, is a common factor of both terms. Therefore the tails of the two members of a pair are $\pm T_X/\v z^2$ exactly, as operators and for every value of $U(\v z)$: the parts linear in $U(\v z)$, the $\v z$-independent parts and the parts with $U(\v z)U^T(\v z)$ cancel together, and nothing has to be assumed about the Wilson line at large distance beyond $|U^{ab}|\le1$; for a target of finite size, $U(\v z)\to1$, the cancellation is a special case. The complex-conjugate half is the complex conjugate of each pair and cancels pair by pair in the same way. Numerically, $\v z^2$ times the sum of the two members decreases like $1/|\v z|$ in every direction, i.e.\ the sum is $O(|\v z|^{-3})$, for all four pairs, $SU(2)$ and $SU(3)$, in the two- and three-$\rho$ sectors, in the colour-singlet projection, with $U(\v z)=1$ and with a random $U(\v z)$ in every direction; the angular mean of $[\kappa_X(w)-\kappa_X(w')]V_XC_X$ falls like $|\v z|^{-4}$.

\paragraph{Kinematic constraint.} The pairing leaves the JIMWLK term unchanged. Equivalently, the soft remainders can be combined with the $\log(k^+/\Lambda)$ terms of the same pairs: with $L=\log(k^+/\Lambda)$,
\begin{align}\label{eq:kinconstr}
&\kappa(w)\int_{\Lambda/k^+}^{1}\frac{d\xi}{\xi}\,e^{-i\xi\,\v k\cdot(\v z-\v w)}-\kappa(w')\int_{\Lambda/k^+}^\infty\frac{d\xi}{\xi(1+\xi)}\,e^{-i\xi\,\v k\cdot(\v z-\v w')}\notag\\
&\quad=L\big[\kappa(w)-\kappa(w')\big]+\kappa(w)\,\mathcal E\big(\v k\cdot(\v z-\v w)\big)-\kappa(w')\,\Psi\big(\v k\cdot(\v z-\v w')\big)+O\Big(\frac{\Lambda}{k^+}\Big)\,.
\end{align}
In this form each kernel of the JIMWLK term that carries a recoil phase comes with its own rapidity interval, $\simeq\log\frac{k^+}{\Lambda|\v k\cdot(\v z-\v w)|}-\gamma_E-\frac{i\pi}{2}{\rm sgn}(\v k\cdot(\v z-\v w))$ for $|\v k\cdot(\v z-\v w)|\gg1$ (resp.\ with $\v w'$) and $\log\frac{k^+}{\Lambda}$ for $|\v k\cdot(\v z-\v w)|\ll1$: the soft gluon has a rapidity below that of the measured gluon, as in the kinematically constrained formulations of forward production~\cite{Altinoluk:2014eka,Ducloue:2016shw} (the constraint of~\cite{Ducloue:2016shw} has not been compared in detail). For pairs~A--C the terms with $L$ are the contributions of these rows to the JIMWLK term; for pair~D only the real half is in the JIMWLK term (see the next subsection).

\subsection{The kernel of Row I.III}\label{sec:fin_GIIII}

The term of $\mathfrak K^{(z)}_{\rm I.III}$, Eq.~\eqref{eq:K-z_rm_I_III}, that multiplies $K(y,w,z)\,P\cdot\Kv(x-w)$ is of a different kind. With $a_t=(x-w)-t(x-z)$,
\begin{align}\label{eq:ell}
\ell(z,w,x)\equiv-\frac12\log\frac{(z-w)^2}{(x-w)^2}=-\int_0^1\frac{dt}{t}\Big[1-\frac{(x-w)\cdot a_t}{a_t^2}\Big]\,,
\end{align}
so that $\ell\,\Kv^i(x-w)$ is the LO kernel minus the LO kernel with the source shifted by the recoil; the soft integral is cut off at $t\simeq|\v x-\v w|/|\v z-\v w|$, without a phase, and nothing of this term is in the JIMWLK list. The term is integrable at $\v z\to\v w$ ($\ell K(y,w,z)\sim\log|\v z-\v w|/|\v z-\v w|$), but at large $\v z$ it has a tail $\propto\ell/\v z^2$ with $\ell\to-\log\frac{|\v z|}{|\v x-\v w|}$, an isotropic tail without the odd part of $\mathcal E_\infty$. In the kernel~\eqref{eq:K-z_rm_I_III} it is therefore written as
\begin{align}
\ell=\Big[\ell-\mathcal E\big(\v k\cdot(\v z-\v w)\big)\Big]+\mathcal E\big(\v k\cdot(\v z-\v w)\big)\,,\qquad
\mathcal E\big(\v k\cdot(\v z-\v w)\big)=\int_0^1d\xi\Big[\frac1\xi\Big]_+e^{-ik\cdot(z-w)\,\xi}\,.
\end{align}
The second part is the Group~I member of pair~D. The first part is combined with the other terms of $\mathfrak K^{(z)}_{\rm I.III}$, the first three groups of terms of Eq.~\eqref{eq:K-z_rm_I_III} (24 terms when written out). Denoting by $\widehat{\mathfrak K}$ the kernel $\mathfrak K^{(z)}_{\rm I.III}$ without its last term, at fixed $\v x,\v y,\v x',\v w'$ and $|\v z|\to\infty$,
\begin{align}\label{eq:Mhatclaim}
\int d^2w\,e^{-i\v k\cdot\v w}\,\widehat{\mathfrak K}=O\Big(\frac{\log|\v z|}{|\v z|^3}\Big)+\big(\text{terms}\propto e^{-i\v k\cdot\v z}\big)\,,
\end{align}
i.e.\ after the $\v w$ integration there is no non-oscillating $1/\v z^2$ term in any direction, for any $U(\v z)$: the colour factor does not depend on $\v w$ and is a spectator. The argument is as follows.
\begin{enumerate}
\item The first three groups of terms of the kernel have the large-$\v z$ form
\begin{align}\label{eq:saw}
\frac{\varphi}{\v z^2}\,\frac{P\cdot e_\perp}{|\v x-\v w|}+O\Big(\frac{\log|\v z|}{|\v z|^3}\Big)\,,\qquad-\pi<\varphi<\pi\,,
\end{align}
where $\varphi$ is the angle from $\v x-\v w$ to $\v z$ and $e_\perp$ the unit vector $(\v x-\v w)/|\v x-\v w|$ rotated by $+\frac\pi2$. It comes from the terms with the arccosine, has zero angular mean, but does not vanish pointwise; Eq.~\eqref{eq:saw} was established numerically.
\item By Eq.~\eqref{eq:recoilfns} and $K(y,w,z)\to1/\v z^2$, the term with $\ell-\mathcal E$ tends to $\frac{P\cdot\Kv(x-w)}{\v z^2}\big[\gamma_E+\log(|\v k||\v x-\v w|)+c(\theta)\big]$, up to terms $\propto e^{-i\v k\cdot(\v z-\v w)}$, with $c(\theta)=\log|\cos\theta|+\frac{i\pi}{2}{\rm sgn}(\cos\theta)$ and $\theta$ the angle between $\v k$ and $\v z$.
\item With $\v r=\v x-\v w$ and $e^{-i\v k\cdot\v w}=e^{-i\v k\cdot\v x}e^{i\v k\cdot\v r}$,
\begin{align}\label{eq:sawid}
\int d^2r\,e^{i\v k\cdot\v r}\,\Kv^i(r)\Big[\gamma_E+\log(|\v k||\v r|)+c(\theta)\Big]&=+\frac{2\pi i\,\hat k^i}{|\v k|}\,\chi(\theta)\,,\notag\\
\int d^2r\,e^{i\v k\cdot\v r}\,\varphi\,\frac{e_\perp^i}{|\v r|}&=-\frac{2\pi i\,\hat k^i}{|\v k|}\,\chi(\theta)\,,\notag\\
\chi(\theta)&\equiv\log|2\cos\theta|+\frac{i\pi}{2}\,{\rm sgn}(\cos\theta)\,,
\end{align}
which follow from the radial integrals $\int_0^\infty dr\,e^{i\kappa r}=i\mathcal P\frac1\kappa+\pi\delta(\kappa)$ and $\int_0^\infty dr\,e^{i\kappa r}[\gamma_E+\log(|\v k|r)]=\frac{i}{\kappa}\big[-\log\frac{|\kappa|}{|\v k|}+\frac{i\pi}{2}{\rm sgn}\,\kappa\big]$, $\kappa=|\v k|\cos\alpha$ ($\alpha$ the angle of $\v r$ with $\v k$; with a convergence factor $e^{-\epsilon r}$), followed by the angular integrals. The two integrals cancel, for every $\theta$. Each of the two Fourier integrals is unique, i.e.\ independent of the regulator and of the order of the radial and angular integrals (Section~\ref{sec:fin_integrable}); what would be only conditionally convergent is the subsequent $\v z$ integral of each term alone, since $\chi$ has zero angular mean. In their sum nothing conditional is left.
\end{enumerate}
Eq.~\eqref{eq:sawid} was checked numerically with the convergence factor $e^{-\epsilon r}$ and with two anisotropic factors (deviations $\le1.2\times10^{-3}$ at $\epsilon=10^{-4}$, falling like $\epsilon$), and Eq.~\eqref{eq:Mhatclaim} with the full kernels: at $|\v z|=2560$ the two parts of $\v z^2\int_we^{-ik\cdot w}\widehat{\mathfrak K}$ are $6.933-2.690i$ and $-6.906+2.690i$, against the prediction $\pm(6.909-2.687i)$, and $\v z^2$ times their sum decreases like $1/|\v z|$ in all directions tested.

The order of the integrations is the definition of the measured momentum (the unmeasured gluon is summed over at fixed $\v k$), not a choice: at fixed $\v w$ the kernel has a $1/\v z^2$ tail whose angular mean does not vanish, so that the other order does not exist; after the $\v w$ integration, which is unique for any regulator of the measured gluon, the tail is absent in every direction and the $\v z$ integral converges absolutely. In summary, the term of Row~I.III with $\ell\,K(y,w,z)\,P\cdot\Kv(x-w)$ is cancelled at large distance (1) in its $\log|\v z|$ part by Row~III.VI, which has the same coefficient and colour structure and the kernel with $\v w\to\v w'$ (pair~D), and (2) in its remaining $1/\v z^2$ part by the tail~\eqref{eq:saw} of the other terms of the same row, after the integration over $\v w$. For pair~D only the real half, from Row~III.VI, is in the JIMWLK term; the virtual half is not, since Row~I.III carries $\ell$ in place of $\log(k^+/\Lambda)$ in this term. The term with the kernel $\mathfrak K^{(w)}_{\rm I.III}$ is cancelled at large distance by Row~III.V (pair~A). At $\v z\to\v w$ both terms are integrable.

\subsection{Instantaneous terms of Group III}\label{sec:fin_inst}

The instantaneous terms of Rows~III.III, III.V and~III.VI, the last term of $\mathfrak K^{(y')}_{\rm III}$ and the kernel $\mathfrak K^{(z)}_{\rm III}$, depend on $\xi$ only through the phase, and have no pointwise limit for $\vee\to\infty$. With $s=\v k\cdot(\v z-\v y')$ and the $\v y'$ integral done first, they are given by the distribution~\eqref{eq:inst_xi}. They are finite. With $\v r=\v y'-\v z$, the combination of the term with $\Kv(x'-y')$ and the Wilson lines at $\v y'$ and of the term with $\Kv(x'-z)$ and the Wilson lines at $\v z$ vanishes at $\v r\to0$, so that $\Kv(y'-z)$ times this combination is bounded there. Integrating first over the component of $\v r$ perpendicular to $\v k$, symmetrically at large distance, gives a function of the component $r_\parallel$ along $\v k$ on which $\pi\delta(s)$ and $\mathcal P\frac1s$ act as on a continuous function, and the phase makes the $r_\parallel$ integral converge at large distance. The two terms must stay together: each alone grows like $\log\vee$, which is the BFKL term. Numerically, with a smooth test function, the longitudinal integral cut at $\vee$ approaches the distributional value like $k^+/\vee$, and the distributional form reproduces the limit to $7\times10^{-8}$. The symmetric cut-off at large distance is sufficient but not necessary: the only part of the perpendicular integral that is not absolutely convergent is independent of $s$ and is annihilated by $\int ds\,e^{-is}[\pi\delta(s)-i\mathcal P\frac1s]=0$, so that any cut-off at large distance gives the same result.

Although the integrand does not decay in $\xi$ at fixed transverse positions, the instantaneous terms do not come from $p^+\sim\vee$, where the eikonal approximation fails. After the $\v y'$ integral the distribution in $\xi$ falls like $(1+\xi)^{-2}$ for Wilson lines that are differentiable, so that $p^+>P$ contributes $O(k^+/P)$, or $O((k^+/P)^\alpha)$ if the Wilson lines are only H\"older continuous with exponent $\alpha$. The bulk comes from $p^+\lesssim k^+\max\{1,1/(|\v k|\ell)\}$, with $\ell$ the shortest transverse scale of the bracket; the eikonal approximation requires $\vee$ to be large compared to this. With a smooth separation $f(p^+/\vee)$ of soft and valence modes ($f(0)=1$) the limit is the same for every $f$, and the whole result depends on $f$ only through $\log\frac{\vee}{k^+}\to\log\frac{\vee}{k^+}+c_f$, $c_f=\int_0^\infty\frac{du}{u}[f(u)-\theta(1-u)]$, in the BFKL term, a change of the rapidity-factorisation scale compensated by the evolution of the projectile. These statements were checked with model brackets that have the short-distance behaviour of the actual one, not with the actual bracket with Wilson lines; that the same $c_f$ appears in every row is argued from the $1/p^+$ form of the integrands, not checked row by row.

\subsection{Conditions and convergence}\label{sec:fin_integrable}

With the pairs of Section~\ref{sec:fin_soft} and the order of integration of Section~\ref{sec:fin_GIIII}, every finite term is integrable over all transverse positions at fixed positions $\v w$, $\v w'$ of the measured gluon, under the following conditions.
\begin{itemize}
\item[(C1)] $\Lambda\to0$ and $\vee\to\infty$ are taken after the subtractions~\eqref{eq:fin_def}, at fixed $\epsilon$ (leading power of the eikonal approximation, Section~\ref{sec:dg_poles}); terms $O(\Lambda/k^+)$ and $O(k^+/\vee)$ are dropped, and the terms $\log(1-k^+/\vee)$ go with the JIMWLK structures of their rows.
\item[(C2)] The identities that relate the colour structures of the two members of a pair use only a relabelling of colour indices and $U(\v z)U^T(\v z)=1$; they are operator identities, valid for any $U$. The statements about the $\v z$ integrand hold at fixed values of all other coordinates after the relabellings of integration variables listed in Table~\ref{tab:pairs}.
\item[(C3)] In Row~I.III the $\v w$ integral, i.e.\ the Fourier transform at the measured momentum, is done before the $\v z$ integral (Section~\ref{sec:fin_GIIII}).
\item[(C4)] Terms $\propto e^{-i\v k\cdot\v z}/\v z^2$ (from the end point $\xi\to1$ in Group~I, and from the regions in which the measured gluon is at $\v z$) are treated as convergent Fourier integrals.
\item[(C5)] No assumption on $U(\v z)$ at large distance is made; the orthogonality of the adjoint Wilson lines is used.
\item[(C6)] The vanishing energy denominators of $\mf B_3$ are principal values.
\item[(C7)] In the instantaneous terms of Group~III the $\v y'$ integral is done before the $\xi$ integral (Section~\ref{sec:fin_inst}).
\item[(C8)] The large-$\v z$ form~\eqref{eq:saw} is established numerically; its Fourier identity~\eqref{eq:sawid} is analytic.
\item[(C9)] The cancellations were verified on a projectile of charges at two sites (all assignments of the charge positions to the sites, including the contact terms of the three-$\rho$ products), with random classical $\rho$ for the three-$\rho$ part, for a general two-$\rho$ coefficient matrix and for its colour-singlet projection.
\item[(C10)] In the region in which both final-state gluons are far from the projectile together ($\v w,\v w',\v z\to\infty$ at fixed $\v w-\v z$), the finite terms of Row~II.IV have the large-distance tail of Section~\ref{sec:dg_tail}. It cancels at the level of the amplitudes against the same region of the genuine two-$\rho$ parts of Row~II.I and of Row~III.III, which therefore have to be integrated together with Row~II.IV under one integral over the position of the pair; the three rows carry the tail with the weights $1\!:\!1\!:\!-2$ (Section~\ref{sec:dg_tail}, Table~\ref{tab:tails}). The region in which only part of the pair is far does not contribute at this order: the three rows are, at fixed $\xi$, $|\mathcal A_B-\mathcal A_{CA}|^2$ with the measured gluon Fourier transformed, and in $\mathcal A_B-\mathcal A_{CA}$ the joint-limit coefficient cancels. This was verified numerically on the full $\v w$ plane for one charge and a target uniform at large distance. For a colour-neutral projectile and a target that is not uniform at large distance the argument is that the non-uniform remainder multiplies $\int_x\Kv(\v x-\v Y)\rho(\v x)=O(\v Y^{-2})$, which was not checked numerically. For a coloured projectile and a target that is not uniform at large distance a remainder proportional to the total colour charge of the projectile is left; it is the non-singlet term of the JIMWLK formula and is absent for a colour-neutral projectile.
\end{itemize}

\paragraph{Convergence.} All integrals over a coordinate that carries the measured phase ($\v w$, $\v w'$, $\v y'$, $\v y-\v y'$, and $\v z$ in the terms $\propto e^{-i\v k\cdot\v z}$) are Fourier integrals of locally integrable, tempered functions whose transforms are continuous at the measured momentum. For such a function $f$ and any regulator $\phi(\epsilon\v r)$ with $\phi(0)=1$ and integrable Fourier transform (an exponential or Gaussian factor, of any direction dependence, or a smooth cut-off of any shape),
\begin{align}\label{eq:FTlemma}
\lim_{\epsilon\to0}\int d^2r\,e^{-i\v k\cdot\v r}\,\phi(\epsilon\v r)\,f(\v r)=\hat f(\v k)\,,
\end{align}
since the left-hand side is the convolution of $\hat f$ with an approximate identity. These integrals, like $\int_we^{-i\v k\cdot\v w}\Kv(x-w)$ at LO, are therefore unique and independent of the order of their radial and angular parts. With them done first, every remaining integral, over the positions of the unmeasured gluon and of the charges and over the longitudinal fractions, converges absolutely in the grouping of Section~\ref{sec:fin_together}, with two exceptions. The first are the principal values $\mathcal P\frac1s$ of the instantaneous terms, which act on H\"older-continuous coefficient functions. The second is the collinear point $\v z\to\v x$ of Row~I.V, where the kernel $\mathfrak K^{\rm quad}_{\rm I.V}$ together with the $\frac{11}{6}$ term behaves as a pure second harmonic in the angle, $a_2\cos2(\varphi-\varphi_0)/|\v z-\v x|^2$ (mean and first harmonic vanish; checked numerically down to $|\v z-\v x|=10^{-4}$): it is conditionally convergent, gives zero with any rotation-invariant cut-off around $\v z=\v x$ and a finite, different value for a non-circular one. The rotation-invariant choice is the one consistent with the $\frac{11}{6}$ subtraction, which comes from the dimensionally regularised (rotation-invariant) integral; an absolutely convergent form is obtained by subtracting $\theta(R_0-|\v z-\v x|)\,a(\varphi)/|\v z-\v x|^2$, which integrates to zero with a rotation-invariant measure. If the transverse indices are continued to $d_\perp=2-\epsilon$, the traceless tensor $\hat r^i\hat r^m-\frac12\delta^{im}$ acquires a trace $O(\epsilon)$, which multiplies the $1/\epsilon$ of the collinear subtraction and could give a finite local term; whether it is present has not been decided (Section~\ref{sec:openpoints}). The two integrals of Eq.~\eqref{eq:sawid} and the two orders of the $\xi$ and $\v y-\v y'$ integrals in Row~II.IV were evaluated with different regulators and orders, with the same result. Hence none of (C3), (C7) and the orders of integration in Eq.~\eqref{eq:sawid} and in Row~II.IV is a prescription that could shift the result by a finite amount. The checks are exact identities, power counting at all coincidences, pointwise tails in all directions and Fourier integrals at fixed values of the other coordinates; a complete numerical integration of the finite terms over all transverse positions with Wilson lines has not been carried out.

\paragraph{Contact terms.} All contact kernels are integrable. Every contact line was power-counted at every coincidence of its integration points (pairs and all larger clusters, in particular the coincidences of the merged charge point with $\v z$, $\v w$, $\v w'$, $\v y'$ and the other charge), with the colour factors as written: each line is absolutely integrable at each coincidence. Only one kernel is not integrable by itself, the instantaneous kernel $\mathfrak K^{(z)}_{\rm III}$ with the charge at $\v x'$ put at $\v x$, which contains $1/(\v x-\v z)^2$; at $\v z\to\v x$ its colour factors vanish identically, so that the line is integrable for any target with H\"older continuous Wilson lines. The kernel $\mathfrak K_{\rm I.II}$ is homogeneous of degree $-1$ in $(\v x-\v w,\v y-\v w)$, bounded at $\v x\to\v y$ and $\propto\log r/r$ at $\v x\to\v w$, and integrable (analytic estimate; its coordinate form is not available). The large-distance tails of the contact lines cancel pair by pair as at separate points. These statements were checked numerically, by fits of the exponents at scales $10^{-2}$ to $10^{-5}$ for three base configurations per line and two smooth, non-uniform $SU(3)$ Wilson-line fields; this is power counting, not a proof for every degenerate configuration.

\subsection{The tail of Row II.IV: term-by-term check}\label{app:tail_check}

The cancellation of the large-distance tail of Row~II.IV against Rows~II.I and~III.III (Section~\ref{sec:dg_tail}) was verified on the finite terms of the three rows. In the joint limit the three rows have the same tail density in $t=p^+/k^+$,
\begin{align}\label{eq:dg_taildens}
\rho(t)=\frac{(t^2+t+1)^2+t^2\cos2\varphi}{t(1+t)^4}\,,
\end{align}
with colour factors $+N_c\Gamma$, $+N_c\Gamma$ and $-N_c\Gamma$ (the last twice, with its complex conjugate), so that the sum vanishes at every $t$. Integrated, in units of $\mathcal U\equiv2N_c\,\mathcal C\frac{4\pi}{\v k^2}\int\frac{d^2Y}{\v Y^2}G_\infty$, the tails of the rows as written are given in Table~\ref{tab:tails}.
\begin{table}[t]\centering
\begin{tabular}{lc}\hline
row & tail$/\,\mathcal U$\\\hline
II.IV: DGLAP $\frac1{12}$, RC $-\frac{11}{12}$, remnant $-1$, tensor $\frac16\cos2\varphi$ & $-\frac{11}{6}+\frac16\cos2\varphi$\\
II.I, second part & $-\frac{11}{6}+\frac16\cos2\varphi$\\
III.III, second part, with c.c. & $+\frac{11}{3}-\frac13\cos2\varphi$\\
II.I and III.III, first parts & $0$\\\hline
sum & $0$\\\hline
\end{tabular}
\caption{Large-distance (soft-parent) tails of the finite terms in the joint limit $\v w,\v w',\v z\to\infty$ at fixed $\v w-\v z$.}\label{tab:tails}
\end{table}
In the first parts the measured gluon on one side is tied to the projectile by the Fourier transform, and there is no tail. In the BFKL and JIMWLK sectors each row carries $(1+1-2)\times$ the same logarithm; the piece moved between sectors by the recoil rewriting of Row~III.III is the companion of its logarithmic term and together with it reproduces the unsubtracted term. The rule $\int d^2z/\v z^2=0$ plays no role here, since the poles it removes multiply $\delta^2(\v k)$. The colour factors were checked with random $SU(2)$ and $SU(3)$ Wilson lines: for Wilson lines that are uniform at large distance the cancellation is exact (to $10^{-14}$), and for a colour-neutral projectile only the terms bilinear in $U(\v x)$, $U(\v x')$ survive, with factors $+N_c,+N_c,-N_c$ independent of the far Wilson lines, so that the tail cancels for any target. For a coloured projectile and a target that is not uniform at large distance a remainder proportional to the total charge survives; it is the non-singlet term of the JIMWLK formula, not a defect of the finite terms. The region in which only part of the pair is far does not contribute at this order (Section~\ref{sec:fin_integrable}, condition (C10)). Not verified are the interchange of the $p^+$ integral with the joint limit in Row~III.III (checked numerically for Row~II.I only) and a full transverse integration with Wilson lines.
\par\endgroup
